% ----------------------------------------------
% Capítol 0
% ----------------------------------------------
\chapter{Chaos}

\epigraf{I la terra estava desordenada i buida, i les tenebres estaven sobre la faç de l'abisme, i l'Esperit de Déu es movia sobre la faç de les aigües.}{Gènesi 1:2}

\lettrine[loversize=0.2, lines=2]{E}{}n un temps llunyà, abans que el coneixement fos reunit sota l’església del OneDrive, els arxius de la saviesa física eren com fragments dispersos en l’infinit, sense ordre ni coherència. Un caos que es mantenia amagat entre els racons d'internet, més enllà de l'abast dels estudiants que, com naus perdudes, lluitaven per trobar el camí entre els arxius perduts. El món de la física, fins a aquell moment, era un conjunt d’intents d'organització fallits, una successió d’emmagatze\-matges caòtics a Google Drive i Mega, llocs plens de carpetes desordenades i arxius amagats en llocs desconeguts.

Aquell temps era com un mar en calma però profundament misteriós, on els arxius existien sense saber que havien de ser alguna cosa més, on els documents eren sols miratges d’allò que es podia aconseguir, sense una idea clara de com donar-li forma al coneixement. Els estudiants, joves i valents, havien de lluitar no només contra els misteris de l'univers, sinó contra el mateix caos que els envoltava: les carpetes dels seus dispositius plenes de noms incomprensibles, els arxius de física desordenats com si fossin estrelles en una nebulosa llunyana. Les lleis que regien aquests sistemes eren tan críptiques com els principis fonamentals de la mecànica quàntica, i qualsevol intent de desxifrar-los semblava tan il·limitat com l’univers mateix.

Sota aquesta capa de desordre i ignorància, hi havia pocs valents que intentaven organitzar aquest caos. A Google Drive, els documents eren nombrosos, però es perdien en un oceà de carpetes que no tenien cap tipus de relació entre elles. Les llegendes parlaven d’arxius oblidats, de fòrums plens de descàrregues en formats incomprensibles, de persones intentant resoldre problemes amb la mateixa desorganització que regnava en els seus arxius.

D’altra banda, el Mega, amb el seu espai infinit, semblava prometre l’ordre, però era només una il·lusió. Les promeses d'una gestió eficaç es van desfer ràpidament davant de la immensitat desordenada que oferia. Els enllaços es perdien com a partícules en un accelerador, i els arxius no trobaven el seu lloc en aquest vast i infinita biblioteca. Els estudiants, atrapats entre aquest caos, lluitaven per salvar-se, però els més savis entenien que la clau no era només en les seves mans, sinó en el poder de la connexió.

Va ser així com, en aquell moment de caos, va sorgir una nova esperança, una promesa d'unitat en el desconegut. Els joves estudiants, que estaven condemnats a sobreviure en aquest món desordenat, buscaven una sortida. Els arxius perduts, les versions de documents incomplertes, les notes de classe disperses, tot esperava ser recol·lectat. Sabien que un dia, algú vindria a posar ordre a aquest caos que governava les seves vides acadèmiques.

La història del OneDrive Física, tal com la coneixem, no havia començat encara. Però els primers passos en aquell viatge ja s’estaven preparant en l’ombra, enmig de la confusió. El caos que regnava seria la llavor de la creació d'una nova era, una època de coneixement compartit, on cada arxiu tindria el seu lloc i cada estudiant el seu accés a l’infinit saber.
